% FILE: 05_evidence_receipts_and_gates.tex
% Project: experiments-visualizer-nextjs
% Role: web-based-pi-visualization-surface

\section{Evidence, Receipts, and Gates}
\label{sec:experiments-visualizer-nextjs-evidence}

\subsection{Tests}
\label{sec:experiments-visualizer-nextjs-evidence-tests}

No tests exist for the visualizer-nextjs project. No testing framework is declared in
\texttt{package.json} --- neither Jest, Vitest, nor Playwright appears as a dependency or
devDependency. No test files are present in the \texttt{src/} directory or at the project
root. The \texttt{npm run lint} command is the only quality-gate command configured.

This absence is a first-class defect under the Van der Aalst Constitution (cf.\
\texttt{{\textasciitilde}/.claude/rules/process-mining-chicago-tdd.md}): \emph{``Model-vs-log
mismatch is not a discrepancy. It is a defect.''} In the context of the visualizer the
analogous principle holds: a CSS class that claims to represent a process-intelligence
state (e.g., \texttt{.alignment-block.sync}) without a component test that verifies it is
applied only in the presence of a genuine sync move is a visual claim without process
evidence. The gap between the design system's class vocabulary and the absence of tests
is precisely such a defect.

The minimum test surface required for ALIVE graduation is:
\begin{enumerate}
  \item Unit tests for each alignment block move type, asserting the correct CSS modifier
        is applied for each of the three lawful move types and no other.
  \item Unit tests for ledger block card states, asserting \texttt{.healthy} and
        \texttt{.tampered} are mutually exclusive and driven by the verification flag.
  \item Integration tests that mount the Petri net panel with a known marking and assert
        the correct set of places bears the \texttt{.active} modifier.
  \item An end-to-end test that loads the dashboard with fixture data and verifies all
        five subsystem panels render without error.
\end{enumerate}

\subsection{Receipts}
\label{sec:experiments-visualizer-nextjs-evidence-receipts}

No receipts exist for the visualizer-nextjs project. The \texttt{.next/BUILD\_ID} file
(\texttt{9NFA-LnYiAijWLGnz8L79}) is a build-system identifier, not a cryptographic receipt
in the sense used by the research program. It does not carry a hash of the build inputs,
a timestamp signed by a process-intelligence operator, or a link to upstream experiment
receipts.

The parent experiments corpus does hold a verified receipt chain (documented in
\texttt{EVIDENCE\_CHAIN\_TRACE.md} and the \texttt{checkpoint\_\_experiments\_complete.md}
file dated 2026-05-31 with an EXPERIMENTS\_COMPLETE verdict). However, this receipt covers
the experiment fixture corpus --- the Rust-typed Raw~$\to$~Parsed~$\to$~Admitted~$\to$~Receipt
evidence chain --- not the web visualization surface. The visualizer-nextjs project is not
covered by any existing receipt.

The minimum receipt required for ALIVE graduation is a project-specific receipt that:
\begin{itemize}
  \item Records the \texttt{BUILD\_ID} of the dashboard build that first rendered
        live process-intelligence data.
  \item Links to the upstream experiment receipt that supplied the data.
  \item Records the hash of the rendered \texttt{index.html} output, confirming the
        page content at the time of receipt issuance.
  \item Is signed by the same receipt-chain operator used by the parent experiments corpus.
\end{itemize}

\subsection{Checkpoints}
\label{sec:experiments-visualizer-nextjs-evidence-checkpoints}

No project-specific checkpoint file exists for the visualizer-nextjs project. The parent
experiments checkpoint (\texttt{checkpoint\_\_experiments\_complete.md}) covers the experiment
fixture corpus and records an EXPERIMENTS\_COMPLETE verdict; it does not cover the web
visualizer. The audit completeness document (\texttt{audit\_\_experiment\_completeness.md})
enumerates what is and is not covered by the parent checkpoint but does not extend coverage
to the visualizer.

For the visualizer to receive its own checkpoint, a designated checkpoint file must be created
in the project directory recording:
\begin{itemize}
  \item The project slug (\texttt{experiments-visualizer-nextjs}).
  \item The verdict (\texttt{ALIVE}, \texttt{PARTIAL}, or \texttt{FAILED}).
  \item References to all receipts and test results that justify the verdict.
  \item The date of the verdict and the signing agent identity.
\end{itemize}

\subsection{ALIVE/PARTIAL/BLOCKED Status}
\label{sec:experiments-visualizer-nextjs-evidence-status}

\textbf{Current status: PARTIAL.}

The following evidence supports the PARTIAL designation:

\begin{itemize}
  \item[\checkmark] The Next.js 16 scaffold compiles successfully (confirmed by
        \texttt{.next/BUILD\_ID} and \texttt{build-diagnostics.json}).
  \item[\checkmark] The CSS design system (\texttt{globals.css}, 973 lines) encodes a
        complete and structurally sound visual vocabulary for all five process-intelligence
        subsystems.
  \item[\checkmark] The project dependency surface (Next.js 16.2.6, React 19.2.4,
        TypeScript 5, Tailwind CSS v4) is correctly declared and resolvable.
  \item[$\times$] \texttt{page.tsx} is the stock welcome page; no process-intelligence
        React components have been implemented.
  \item[$\times$] No data-fetching layer or API routes are present.
  \item[$\times$] No tests exist.
  \item[$\times$] No project-specific receipt exists.
  \item[$\times$] No project-specific checkpoint exists.
\end{itemize}

Under the research program's ALIVE definition, all five of the missing items must be
remediated before an ALIVE verdict can be issued. The project may not be declared ALIVE
on the basis of the CSS design system alone, regardless of its completeness and quality.
